% !TEX program = xelatex
% !TEX options = -synctex=1 -interaction=nonstopmode -file-line-error
% example-abstract.tex
% Demonstrates rossabstract: conference header, structured abstract
% with all four sections, and keywords.

\documentclass[palette=lancet, wordlimit=300]{rossabstract}

% ── Metadata ─────────────────────────────────────────────────────
\title{Slow-Wave Sleep Predicts Delayed Memory via Depressive Symptoms
       in Community-Dwelling Older Adults}
\author{Ross A.\ Dunne\textsuperscript{1},
        Sarah K.\ Mitchell\textsuperscript{2},
        Ann-Marie Moriarty\textsuperscript{1}}
\affiliation{%
  \textsuperscript{1}Department of Psychology, University of Limerick;
  \textsuperscript{2}School of Medicine, Trinity College Dublin}
\correspondingauthor{ross.dunne@ul.ie}
\conference{53rd Annual Meeting of the International Neuropsychological
            Society --- New York, NY, USA --- February 2026}
\abstracttype{Oral Presentation}
\keywords{sleep architecture, slow-wave sleep, cognition, mediation,
          depressive symptoms, ageing}
\conflictofinterest{None declared.}
\funding{Irish Research Council (GOIPG/2024/1234).}

\begin{document}

\maketitle

\structuredabstract
  % ── Background ─────────────────────────────────────────────────
  {Sleep disturbance is a recognised risk factor for cognitive decline in
   older adults, yet the pathways linking specific sleep-architecture
   parameters to domain-specific cognitive outcomes remain poorly
   understood. Depressive symptoms frequently co-occur with disrupted
   sleep and may act as a mediator of cognitive change.}
  % ── Methods ────────────────────────────────────────────────────
  {Community-dwelling adults aged 60--85 ($N = 312$; 58\% female) completed
   single-night in-lab polysomnography, the DASS-21, and the Repeatable
   Battery for the Assessment of Neuropsychological Status (RBANS).
   Parallel mediation models tested indirect effects of slow-wave sleep
   percentage (SWS\%) on each RBANS index through depressive and anxiety
   symptom severity, controlling for age, sex, and education.}
  % ── Results ────────────────────────────────────────────────────
  {SWS\% was significantly associated with delayed memory
   ($\beta = -0.24$, $p = .001$). The indirect effect through depressive
   symptoms was significant ($\beta = -0.18$, 95\% CI [$-$0.34, $-$0.08]).
   No significant indirect effect was observed through anxiety symptoms
   ($p = .41$). Effects on immediate memory, attention, language, and
   visuospatial domains were non-significant after correction for
   multiple comparisons.}
  % ── Conclusions ────────────────────────────────────────────────
  {Depressive symptoms partially mediate the relationship between reduced
   slow-wave sleep and poorer delayed memory in older adults. Interventions
   targeting mood disturbance in individuals with sleep complaints may
   confer downstream cognitive benefits. Longitudinal replication with
   multi-night polysomnography is warranted.}

\end{document}
